%\begin{center}
%\large \bf \runtitulo
%\end{center}
%\vspace{1cm}
\chapter*{\runtitulo}

\noindent \texit{Java Typestate Checker} es una herramienta de verificación estática de protocolos de objetos. A través de la abstracción de los \textit{typestates}, representando el estado de un objeto a lo largo de la ejecución de un programa, asegura que los métodos de un objeto sean llamados en su orden correcto. Esto mitiga errores y comportamientos inesperados a la hora de usar clases con protocolo en tiempo de compilado. A su vez, la herramienta aporta nuevas anotaciones para especificar otros comportamientos esperados de los programas, tales como qué variables pueden tener el valor \texttt{null} o cuál es el \textit{typestate} esperado para parámetros y valores de retorno en métodos. Este trabajo recopila información sobre sus características e implementación, y la evalúa experimentalmente para entender sus capacidades y límites, así como los de la verificación de protocolos de objetos mediante \textit{typestates} en general.

\bigskip

\noindent\textbf{Palabras claves:} typestate, protocolos de objetos, verificación estática, Java, Java Typestate Checker, orden de invocación de métodos.